\documentclass[12pt]{article}

% Idioma y codificación
\usepackage[spanish, es-tabla, es-notilde]{babel}       %es-tabla para que se titule "Tabla"
\usepackage[utf8]{inputenc}

% Márgenes
\usepackage[a4paper,top=3cm,bottom=2.5cm,left=3cm,right=3cm]{geometry}

% Comentarios de bloque
\usepackage{verbatim}

% Paquetes de links
\usepackage[hidelinks]{hyperref}    % Permite enlaces
\usepackage{url}                    % redirecciona a la web

% Más opciones para enumeraciones
\usepackage{enumitem}

% Personalizar la portada
\usepackage{titling}

% Paquetes de tablas
\usepackage{multirow}

% Para añadir el símbolo de euro
\usepackage{eurosym}

% Paquete para subrayar con diferentes estilos
\usepackage[normalem]{ulem}


%------------------------------------------------------------------------

%Paquetes de figuras
\usepackage{caption}
\usepackage{subcaption} % Figuras al lado de otras
\usepackage{float}      % Poner figuras en el sitio indicado H.


% Paquetes de imágenes
\usepackage{graphicx}       % Paquete para añadir imágenes
\usepackage{transparent}    % Para manejar la opacidad de las figuras

% Paquete para usar colores
\usepackage[dvipsnames, table, xcdraw]{xcolor}
\usepackage{pagecolor}      % Para cambiar el color de la página

% Habilita tamaños de fuente mayores
\usepackage{fix-cm}

% Para los gráficos
\usepackage{tikz}
\usepackage{forest}

% Para poder situar los nodos en los grafos
\usetikzlibrary{positioning}

% Para evitar problemas entre babel y tikz
\usetikzlibrary{babel}


%------------------------------------------------------------------------

% Paquetes de matemáticas
\usepackage{mathtools, amsfonts, amssymb, mathrsfs}
\usepackage[makeroom]{cancel}     % Simplificar tachando
\usepackage{polynom}    % Divisiones y Ruffini
\usepackage{units} % Para poner fracciones diagonales con \nicefrac

\usepackage{pgfplots}   %Representar funciones
\pgfplotsset{compat=1.18}  % Versión 1.18

\usepackage{tikz-cd}    % Para usar diagramas de composiciones
\usetikzlibrary{calc}   % Para usar cálculo de coordenadas en tikz

%Definición de teoremas, etc.
\usepackage{amsthm}
%\swapnumbers   % Intercambia la posición del texto y de la numeración

\theoremstyle{plain}

\makeatletter
\@ifclassloaded{article}{
  % --- Configuración para ARTICLE ---
  \newtheorem{teo}{Teorema}[section]
  \theoremstyle{definition}
  \newtheorem{definicion}{Definición}[section]
  \newtheorem{ejercicio}{Ejercicio}[section]

  \renewcommand{\theteo}{%
    \ifnum\value{section}=0 \arabic{teo}\else \thesection.\arabic{teo}\fi
  }
  \renewcommand{\thedefinicion}{%
    \ifnum\value{section}=0 \arabic{definicion}\else \thesection.\arabic{definicion}\fi
  }
  \renewcommand{\theejercicio}{%
    \ifnum\value{section}=0 \arabic{ejercicio}\else \thesection.\arabic{ejercicio}\fi
  }
}{
  % --- Configuración para BOOK / REPORT ---
  \newtheorem{teo}{Teorema}[chapter]
  \theoremstyle{definition}
  \newtheorem{definicion}{Definición}[chapter]
  \newtheorem{ejercicio}{Ejercicio}[section]

  % Si no hay sección: Capítulo.Teorema (ej. 2.1)
  % Si hay sección:    Capítulo.Sección.Teorema (ej. 2.3.1)
  \renewcommand{\theteo}{%
    \ifnum\value{section}=0
      \thechapter.\arabic{teo}%
    \else
      \thesection.\arabic{teo}% (\thesection ya contiene el número de capítulo)
    \fi
  }
  \renewcommand{\thedefinicion}{%
    \ifnum\value{section}=0
      \thechapter.\arabic{definicion}%
    \else
      \thesection.\arabic{definicion}%
    \fi
  }
  \renewcommand{\theejercicio}{%
    \ifnum\value{section}=0
      \arabic{ejercicio}%
    \else
      \thesection.\arabic{ejercicio}% Formato por defecto: ej. 2.2.1
    \fi
  }
}
\makeatother

% Teoremas dependientes o con contador compartido
\newtheorem{coro}{Corolario}[teo]
\newtheorem{prop}[teo]{Proposición}
\newtheorem{lema}[teo]{Lema}
\newtheorem*{lema*}{Lema}

\theoremstyle{remark}
\newtheorem*{observacion}{Observación}

\theoremstyle{definition}
\newtheorem*{notacion}{Notación}
\newtheorem*{ejemplo}{Ejemplo}
\newtheorem*{ejercicio*}{Ejercicio}


% \renewcommand\qedsymbol{$\blacksquare$}         % Cambiar símbolo QED
%------------------------------------------------------------------------

% Paquetes para encabezados
\usepackage{fancyhdr}
\pagestyle{fancy}
\fancyhf{}

\newcommand{\helv}{ % Modificación tamaño de letra
\fontfamily{}\fontsize{12}{12}\selectfont}
\setlength{\headheight}{15pt} % Amplía el tamaño del índice


%\usepackage{lastpage}   % Referenciar última pag   \pageref{LastPage}
%\fancyfoot[C]{%
%  \begin{minipage}{\textwidth}
%    \centering
%    ~\\
%    \thepage\\
%    \href{https://losdeldgiim.github.io/}{\texttt{\footnotesize losdeldgiim.github.io}}
%  \end{minipage}
%}
\fancyfoot[C]{\thepage}
\fancyfoot[R]{\href{https://losdeldgiim.github.io/}{\texttt{\footnotesize losdeldgiim.github.io}}}

%------------------------------------------------------------------------

% Conseguir que no ponga "Capítulo 1". Sino solo "1."
\makeatletter
\@ifclassloaded{book}{
  \renewcommand{\chaptermark}[1]{\markboth{\thechapter.\ #1}{}} % En el encabezado
    
  \renewcommand{\@makechapterhead}[1]{%
  \vspace*{50\p@}%
  {\parindent \z@ \raggedright \normalfont
    \ifnum \c@secnumdepth >\m@ne
      \huge\bfseries \thechapter.\hspace{1em}\ignorespaces
    \fi
    \interlinepenalty\@M
    \Huge \bfseries #1\par\nobreak
    \vskip 40\p@
  }}
}
\makeatother

%------------------------------------------------------------------------
% Paquetes de cógido
\usepackage{minted}
\renewcommand\listingscaption{Código fuente}

\usepackage{fancyvrb}
% Personaliza el tamaño de los números de línea
\renewcommand{\theFancyVerbLine}{\small\arabic{FancyVerbLine}}

% Estilo para C++
\newminted{cpp}{
    frame=lines,
    framesep=2mm,
    baselinestretch=1.2,
    linenos,
    escapeinside=||
}

% para minted
\definecolor{LightGray}{rgb}{0.95,0.95,0.92}
\setminted{
    linenos=true,
    stepnumber=5,
    numberfirstline=true,
    autogobble,
    breaklines=true,
    breakautoindent=true,
    breaksymbolleft=,
    breaksymbolright=,
    breaksymbolindentleft=0pt,
    breaksymbolindentright=0pt,
    breaksymbolsepleft=0pt,
    breaksymbolsepright=0pt,
    fontsize=\footnotesize,
    bgcolor=LightGray,
    numbersep=10pt
}


\usepackage{listings} % Para incluir código desde un archivo

\renewcommand\lstlistingname{Código Fuente}
\renewcommand\lstlistlistingname{Índice de Códigos Fuente}

% Definir colores
\definecolor{vscodepurple}{rgb}{0.5,0,0.5}
\definecolor{vscodeblue}{rgb}{0,0,0.8}
\definecolor{vscodegreen}{rgb}{0,0.5,0}
\definecolor{vscodegray}{rgb}{0.5,0.5,0.5}
\definecolor{vscodebackground}{rgb}{0.97,0.97,0.97}
\definecolor{vscodelightgray}{rgb}{0.9,0.9,0.9}

% Configuración para el estilo de C similar a VSCode
\lstdefinestyle{vscode_C}{
  backgroundcolor=\color{vscodebackground},
  commentstyle=\color{vscodegreen},
  keywordstyle=\color{vscodeblue},
  numberstyle=\tiny\color{vscodegray},
  stringstyle=\color{vscodepurple},
  basicstyle=\scriptsize\ttfamily,
  breakatwhitespace=false,
  breaklines=true,
  captionpos=b,
  keepspaces=true,
  numbers=left,
  numbersep=5pt,
  showspaces=false,
  showstringspaces=false,
  showtabs=false,
  tabsize=2,
  frame=tb,
  framerule=0pt,
  aboveskip=10pt,
  belowskip=10pt,
  xleftmargin=10pt,
  xrightmargin=10pt,
  framexleftmargin=10pt,
  framexrightmargin=10pt,
  framesep=0pt,
  rulecolor=\color{vscodelightgray},
  backgroundcolor=\color{vscodebackground},
}

%------------------------------------------------------------------------
% La etiqueda IA sirve para indicar a partir de donde le pasamos el preambulo a la IA para el escaneo
% IA

% Comandos definidos
\newcommand{\bb}[1]{\mathbb{#1}}
\newcommand{\cc}[1]{\mathcal{#1}}

% I prefer the slanted \leq
\let\oldleq\leq % save them in case they're every wanted
\let\oldgeq\geq
\renewcommand{\leq}{\leqslant}
\renewcommand{\geq}{\geqslant}

% Si y solo si
\newcommand{\sii}{\iff}

% MCD y MCM
\DeclareMathOperator{\mcd}{mcd}
\DeclareMathOperator{\mcm}{mcm}

% Signo
\DeclareMathOperator{\sgn}{sgn}

% Imagen
\DeclareMathOperator{\Ima}{Im}

% Aplicación identidad
\newcommand{\id}{\mathrm{Id}}

% Letras griegas
\newcommand{\eps}{\epsilon}
\newcommand{\veps}{\varepsilon}
\newcommand{\lm}{\lambda}

\newcommand{\ol}{\overline}
\newcommand{\ul}{\underline}
\newcommand{\wt}{\widetilde}
\newcommand{\wh}{\widehat}

\let\oldvec\vec
\renewcommand{\vec}{\overrightarrow}

% Derivadas parciales
\newcommand{\del}[2]{\frac{\partial #1}{\partial #2}}
\newcommand{\Del}[3]{\frac{\partial^{#1} #2}{\partial #3^{#1}}}
\newcommand{\deld}[2]{\dfrac{\partial #1}{\partial #2}}
\newcommand{\Deld}[3]{\dfrac{\partial^{#1} #2}{\partial #3^{#1}}}


\newcommand{\AstIg}{\stackrel{(\ast)}{=}}
\newcommand{\Hop}{\stackrel{L'H\hat{o}pital}{=}}

\newcommand{\red}[1]{{\color{red}#1}} % Para integrales, destacar los cambios.

% Método de integración
\newcommand{\MetInt}[2]{
    \left[\begin{array}{c}
        #1 \\ #2
    \end{array}\right]
}

% Declarar aplicaciones
% 1. Nombre aplicación
% 2. Dominio
% 3. Codominio
% 4. Variable
% 5. Imagen de la variable
\newcommand{\Func}[5]{
    \begin{equation*}
        \begin{array}{rrll}
            \displaystyle #1:& \displaystyle  #2 & \longrightarrow & \displaystyle  #3\\
               & \displaystyle  #4 & \longmapsto & \displaystyle  #5
        \end{array}
    \end{equation*}
}

%------------------------------------------------------------------------


\newcommand{\myhref}[2]{\href{#1}{\color{blue}\uline{#2}}}

\begin{document}

    % 1. Foto de fondo
    % 2. Título
    % 3. Encabezado Izquierdo
    % 4. Color de fondo
    % 5. Coord x del titulo
    % 6. Coord y del titulo
    % 7. Fecha

    
    % 1. Foto de fondo
% 2. Título
% 3. Encabezado Izquierdo
% 4. Color de fondo
% 5. Coord x del titulo
% 6. Coord y del titulo
% 7. Fecha
% 8. Autor

\newcommand{\portada}[8]{

    \portadaBase{#1}{#2}{#3}{#4}{#5}{#6}{#7}{#8}
    \portadaBook{#1}{#2}{#3}{#4}{#5}{#6}{#7}{#8}
}

\newcommand{\portadaExamen}[8]{

    \portadaBase{#1}{#2}{#3}{#4}{#5}{#6}{#7}{#8}
    \portadaArticle{#1}{#2}{#3}{#4}{#5}{#6}{#7}{#8}
}




\newcommand{\portadaBase}[8]{

    % Tiene la portada principal y la licencia Creative Commons
    
    % 1. Foto de fondo
    % 2. Título
    % 3. Encabezado Izquierdo
    % 4. Color de fondo
    % 5. Coord x del titulo
    % 6. Coord y del titulo
    % 7. Fecha
    % 8. Autor
    
    
    \thispagestyle{empty}               % Sin encabezado ni pie de página
    \newgeometry{margin=0cm}        % Márgenes nulos para la primera página
    
    
    % Encabezado
    \fancyhead[L]{\helv #3}
    \fancyhead[R]{\helv \nouppercase{\leftmark}}
    
    
    \pagecolor{#4}        % Color de fondo para la portada
    
    \begin{figure}[p]
        \centering
        \transparent{0.3}           % Opacidad del 30% para la imagen
        
        \includegraphics[width=\paperwidth, keepaspectratio]{../../_assets/#1}
    
        \begin{tikzpicture}[remember picture, overlay]
            \node[anchor=north west, text=white, opacity=1, font=\fontsize{60}{90}\selectfont\bfseries\sffamily, align=left] at (#5, #6) {#2};
            
            \node[anchor=south east, text=white, opacity=1, font=\fontsize{12}{18}\selectfont\sffamily, align=right] at (9.7, 3) {\href{https://losdeldgiim.github.io/}{\textbf{Los Del DGIIM}, \texttt{\footnotesize losdeldgiim.github.io}}};
            
            \node[anchor=south east, text=white, opacity=1, font=\fontsize{12}{15}\selectfont\sffamily, align=right] at (9.7, 1.8) {Doble Grado en Ingeniería Informática y Matemáticas\\Universidad de Granada};
        \end{tikzpicture}
    \end{figure}
    
    
    \restoregeometry        % Restaurar márgenes normales para las páginas subsiguientes
    \pagecolor{white}       % Restaurar el color de página
    
    
    \newpage
    \thispagestyle{empty}               % Sin encabezado ni pie de página
    \begin{tikzpicture}[remember picture, overlay]
        \node[anchor=south west, inner sep=3cm] at (current page.south west) {
            \begin{minipage}{0.5\paperwidth}
                \href{https://creativecommons.org/licenses/by-nc-nd/4.0/}{
                    \includegraphics[height=2cm]{../../_assets/Licencia.png}
                }\vspace{1cm}\\
                Esta obra está bajo una
                \href{https://creativecommons.org/licenses/by-nc-nd/4.0/}{
                    Licencia Creative Commons Atribución-NoComercial-SinDerivadas 4.0 Internacional (CC BY-NC-ND 4.0).
                }\\
    
                Eres libre de compartir y redistribuir el contenido de esta obra en cualquier medio o formato, siempre y cuando des el crédito adecuado a los autores originales y no persigas fines comerciales. 
            \end{minipage}
        };
    \end{tikzpicture}
    
    
    
    % 1. Foto de fondo
    % 2. Título
    % 3. Encabezado Izquierdo
    % 4. Color de fondo
    % 5. Coord x del titulo
    % 6. Coord y del titulo
    % 7. Fecha
    % 8. Autor


}


\newcommand{\portadaBook}[8]{

    % 1. Foto de fondo
    % 2. Título
    % 3. Encabezado Izquierdo
    % 4. Color de fondo
    % 5. Coord x del titulo
    % 6. Coord y del titulo
    % 7. Fecha
    % 8. Autor

    % Personaliza el formato del título
    \pretitle{\begin{center}\bfseries\fontsize{42}{56}\selectfont}
    \posttitle{\par\end{center}\vspace{2em}}
    
    % Personaliza el formato del autor
    \preauthor{\begin{center}\Large}
    \postauthor{\par\end{center}\vfill}
    
    % Personaliza el formato de la fecha
    \predate{\begin{center}\huge}
    \postdate{\par\end{center}\vspace{2em}}
    
    \title{#2}
    \author{\href{https://losdeldgiim.github.io/}{Los Del DGIIM, \texttt{\large losdeldgiim.github.io}}
    \\ \vspace{0.5cm}#8}
    \date{Granada, #7}
    \maketitle
    
    \tableofcontents
}




\newcommand{\portadaArticle}[8]{

    % 1. Foto de fondo
    % 2. Título
    % 3. Encabezado Izquierdo
    % 4. Color de fondo
    % 5. Coord x del titulo
    % 6. Coord y del titulo
    % 7. Fecha
    % 8. Autor

    % Personaliza el formato del título
    \pretitle{\begin{center}\bfseries\fontsize{42}{56}\selectfont}
    \posttitle{\par\end{center}\vspace{2em}}
    
    % Personaliza el formato del autor
    \preauthor{\begin{center}\Large}
    \postauthor{\par\end{center}\vspace{3em}}
    
    % Personaliza el formato de la fecha
    \predate{\begin{center}\huge}
    \postdate{\par\end{center}\vspace{5em}}
    
    \title{#2}
    \author{\href{https://losdeldgiim.github.io/}{Los Del DGIIM, \texttt{\large losdeldgiim.github.io}}
    \\ \vspace{0.5cm}#8}
    \date{Granada, #7}
    \thispagestyle{empty}               % Sin encabezado ni pie de página
    \maketitle
    \vfill
}
    \portadaExamen{ffccA4.jpg}{Topología I\\Examen XVIII}{Topología I. Examen XVIII}{MidnightBlue}{-8}{28}{2026}{Lucía Benjumeda Maldonado}

    \begin{description}
        \item[Asignatura] Topología I.
        \item[Curso Académico] 2025-26.
        \item[Grado] Doble Grado en Ingeniería Informática y Matemáticas.
        \item[Grupo] Único.
        \item[Profesor] Antonio Alarcón López.
        \item[Descripción] Parcial 2.
        \item[Fecha] 12 de noviembre de 2025.
        \item[Duración] 120 minutos.
    
    \end{description}
    \newpage

    \begin{ejercicio}
    (2,5 puntos) Define espacio topológico T1 y demuestra que $(X, T)$ es T1 si y sólo si $\{x\}$ es cerrado para todo $x \in X$.
    \end{ejercicio}

    \begin{ejercicio}
    Sea $X$ un conjunto y $A \subset X$ con $A \neq \varnothing$ y $A \neq X$. Se considera la siguiente familia de subconjuntos de $X$:
    \[
        B = \{\{x\} \cup A : x \in X\}.
    \]
    \begin{enumerate}
        \item (1 punto) Prueba que existe una única topología $T$ en $X$ para la cual $B$ es base.
        \item (0,5 puntos) Demuestra que si $U \subset X$ con $U \neq \varnothing$, entonces $U \in T$ si y sólo si $A \subset U$.
        \item (1 punto) Identifica las topologías inducidas sobre $A$ y $X \setminus A$.
        \item (1 punto) Estudia si $(X, T)$ es 1AN y 2AN.
    \end{enumerate}
    \end{ejercicio}

    \begin{ejercicio}
    En $\mathbb{R}^2$ denotamos por $d$ a la distancia usual y definimos la familia de subconjuntos $T$ de la siguiente forma:
    \[
        T = \{U \subset \mathbb{R}^2 : (0,0) \notin U\} \cup \{U \subset \mathbb{R}^2 : \exists \varepsilon > 0 \text{ con } B_d((0,0), \varepsilon) \subset U\}.
    \]
    \begin{enumerate}
        \item (1,5 punto) Demuestra que $T$ es una topología en $\mathbb{R}^2$. Compárala con la topología usual de $\mathbb{R}^2$.
        \item (1 punto) Calcula una base de entornos de $(x,y) \in \mathbb{R}^2$ en $(\mathbb{R}^2, T)$.
        \item (1 punto) Calcula la adherencia y el interior de $\mathbb{Q} \times \{0\}$ y $\mathbb{Q} \times (0,1)$ en $(\mathbb{R}^2, T)$.
    \end{enumerate}
    \end{ejercicio}

    \newpage
    \setcounter{ejercicio}{0} % Reiniciar contador de ejercicios
    \noindent

    \begin{ejercicio} 
    (2,5 puntos) Define espacio topológico T1 y demuestra que $(X, T)$ es T1 si y sólo si $\{x\}$ es cerrado para todo $x \in X$.
    
    Un espacio topológico $(X, T)$ es T1 si para cada par de puntos $x, y\in X$, $x\neq y$ existen un $U, V\in T$ con $x\in U$ y $y\in V$ tales que $y\notin U$ y $x\notin V$.\\

    Veamos ahora por doble implicación que $(X, T)$ es T1 si y sólo si $\{x\}$ es cerrado para todo $x \in X$:
    \begin{itemize}
        \item[$\Rightarrow$] Supongamos que $(X, T)$ es T1, vamos a ver que sea $x\in X$ arbitrario entonces $\{x\} = \overline{\{x\}}$. Está claro que $\{x\} \subset \overline{\{x\}}$. Para la otra inclusión sea $y \in \overline{\{x\}}$, entonces para todo $V\in T$ con $y\in V$ se tiene que $V\cap \{x\} \neq \varnothing$, lo que implica que $x\in V$. Si fuese $x \neq y$, como $(X, T)$ es T1 existiría $V' \in T$ con $y \in V'$ y $x \notin V'$, contradicción. Por tanto $x = y$ (y en particular $y \in \{x\}$).
        \item[$\Leftarrow$] Supongamos que $\{x\}$ es cerrado para todo $x \in X$, vamos a ver que $(X, T)$ es T1. Sea $x, y \in X$ con $x \neq y$. Entonces $x \in X \setminus \{y\}$ y $y \in X \setminus \{x\}$. Como $U = X \setminus \{y\}$ y $V =X \setminus \{x\}$ son abiertos (complementario de cerrado), existen $U, V \in T$ tales que $x \in U$ y $y \in V$ y además, $y \notin U$ y $x \notin V$.
    \end{itemize} 
    \end{ejercicio}

    \begin{ejercicio}
    Sea $X$ un conjunto y $A \subset X$ con $A \neq \varnothing$ y $A \neq X$. Se considera la siguiente familia de subconjuntos de $X$:
    \[
        B = \{\{x\} \cup A : x \in X\}.
    \]
    \begin{enumerate}
        \item (1 punto) Prueba que existe una única topología $T$ en $X$ para la cual $B$ es base.\\
        Vamos a probar que $\cc{B} = \{B\subseteq X: B = \{x\}\cup A \text{ con } x \in X\}\subseteq \cc{P}(X)$ es una base de una topología en $X$. Para ello probemos que cumple las siguientes propiedades:
        \begin{enumerate}
            \item[(B1)] $X = \bigcup_{B \in B} B$
            \begin{itemize}
                \item[$\subseteq$] Trivial porque sea $x \in X$, entonces existe un $B\in \cc{B}$ tal que $x \in B = \{x\} \cup A$.
                \item[$\supseteq$] Sea $x \in \bigcup_{B \in B} B$, entonces existe un $B_x\in \cc{B}$ tal que $x \in B_x = \{x\} \cup A$ por definición con $x\in X$.
            \end{itemize}
            \item[(B2)] Sea $B_1, B_2 \in \cc{B}$ y $x \in B_1 \cap B_2$, existe $B_3 \in \cc{B}$ tal que $x \in B_3 \subset B_1 \cap B_2$. Sean $B_1 = \{x_1\}\cup A$ y $B_2 = \{x_2\}\cup A$, distinguimos los siguientes casos:
            \begin{enumerate}
                \item Si $x\in A$ bastaría tomar $B_3 = \{x\} \cup A = A$ porque claramente $x\in B_3 = A \subset B_1 \cap B_2$.
                \item Si $x\notin A$ entonces $x = x_1$ y $x = x_2$ para estar en la intersección. Supongamos que $x = x_1$, entonces tomando $B_3 = B_1 = \{x\} \cup A$ cumple que $x \in B_3 \subset B_1 \cap B_2$.
            \end{enumerate}
        \end{enumerate}
        Así quedaría demostrado que existe una única topología $T$ en $X$ para la cual $B$ es base.
        \item (0,5 puntos) Demuestra que si $U \subset X$ con $U \neq \varnothing$, entonces $U \in T$ si y sólo si $A \subset U$.\\
        Demostremos la doble implicación:
        \begin{itemize}
            \item[$\Rightarrow$] Supongamos que $U \in T$, entonces $U = \bigcup_{B_i\in \cc{B}} B_i = \bigcup_{i\in I} (\{x_i\} \cup A) = (\bigcup_{i\in I} \{x_i\}) \cup A$, por lo que, $A \subset U$.
            \item[$\Leftarrow$] Supongamos que $A \subset U$, entonces para todo $a \in A$ existe un $B_a = \{a\} \cup A = A$ con $a\in B_a\subset A\subset U$ luego $U\in \cc{N}_a, \forall a \in A$. Por otro lado si $b\in U\setminus A$ existe un $B_b = \{b\} \cup A$ con $b\in B_b\subset U$ y por lo tanto $U \in \cc{N}_b$. Luego $U\in \cc{N}_x, \forall x \in U$ lo que implica que $U \in T$.
        \end{itemize}
        \item (1 punto) Identifica las topologías inducidas sobre $A$ y $X \setminus A$.
        
        Comencemos con la topología inducida $(A, T_{|A})$ que sabemos que tiene como base de la topología $\cc{B}_{|A} = \{B\cap A : B \in \cc{B}\} = \{(\{x\} \cup A)\cap A : x\in X\} = \{\varnothing\} \cup\{A\}$. De esta manera si escribimos todo abierto como unicón de sus bases tenemos sea $V\in T_{|A}$ entonces: 
        \begin{itemize}
            \item $V = \bigcup_{B\in\cc{B}_{|A}} B = A$
            \item $V = \varnothing$
        \end{itemize}
        Luego $T_{|A} = \{\varnothing, A\}$ es la topología trivial.\\

        Por otro lado, la topología inducida $(X\setminus A, T_{|X\setminus A})$ tiene como base de la topología $\cc{B}_{|X\setminus A} = \{B\cap (X\setminus A) : B \in \cc{B}\} = \{(\{x\} \cup A)\cap (X\setminus A) : x\in X\} = \{\{x\} : x\in X\setminus A\}$. De esta manera si escribimos todo abierto como unión de sus bases tenemos sea $V\in T_{|X\setminus A}$ entonces $V = \bigcup_{B\in\cc{B}_{|X\setminus A}} B = \bigcup_{x\in V} \{x\} = V$. Con lo que claramente se ve que $T_{|X\setminus A} = \cc{P}(X\setminus A)$ es la topología discreta.
        
        \item (1 punto) Estudia si $(X, T)$ es 1AN y 2AN.\\
        Vamos a probar que $(X, T)$ es 1AN. Para ello observemos que sea $x\in X$ entonces $\cc{B}_x=\{\{x\}\cup A\}$ es una base de entornos de $x$ pues sea $U\in T$ arbitrario, con $x\in U$, como $U \neq \varnothing$ y $U \in T$, por el apartado 2 se tiene que $A \subset U$. Como además $x \in U$, de alguna de las maneras, $\exists V = \{x\} \cup A$ entorno básico de $x$ tal que $x\in V\subset U$.\\
        Habiendo demostrado que en efecto $\cc{B}_x=\{\{x\}\cup A\}$ es base de entornos de $x$, el resultado es inmediato pues $\cc{B}_x$ es finita, de hecho tiene un único elemento, luego es numerable y $(X, T)$ es 1AN.\\
        
        Ahora vamos a demostrar que $(X, T)$ puede ser o no 2AN dependiendo de $A$:
        \begin{itemize}
            \item Si $X\setminus A$ es numerable (o finito), entonces existe una cantidad numerable de puntos $\{x_n\}_{n\in I} = X\setminus A$ (con $I \subseteq \bb{N}$) y por lo tanto tomando $\cc{B}_n = \{A\} \cup \{\{x_n\}\cup A : n\in I\}$ tenemos una familia numerable. Comprobemos que es efectivamente una base de $T$:
            
            Sabemos por el apartado 1 que $\cc{B} = \{\{x\}\cup A : x\in X\}$ es una base de $T$. Basta observar que $\cc{B}_n = \cc{B}$, ya que para todo $x \in X$:
            \begin{itemize}
                \item Si $x \in A$, entonces $\{x\} \cup A = A \in \cc{B}_n$.
                \item Si $x \in X \setminus A$, entonces $x = x_n$ para algún $n \in I$, luego $\{x\} \cup A = \{x_n\} \cup A \in \cc{B}_n$.
            \end{itemize}
            Por tanto, $\cc{B}_n$ es una base numerable de $T$ y $(X, T)$ es 2AN.

            \item Si $X\setminus A$ es no numerable, supongamos por reducción al absurdo que $(X, T)$ tiene una base numerable $\cc{B}_n$. Tomemos $I = X \setminus A$ (no numerable) y la familia de abiertos $\{U_i\}_{i\in I}$ definida por $U_i = \{x_i\} \cup A$ para cada $x_i \in X\setminus A$. Por definición de base, para cada $i\in I$ existe $B_i \in \cc{B}_n$ tal que $x_i \in B_i \subset U_i = \{x_i\}\cup A$. Como $B_i \in T$ y $B_i \neq \varnothing$, por el apartado 2 se tiene que $A \subset B_i$, lo que fuerza a que $B_i = \{x_i\}\cup A$. Si tomamos dos puntos $x_i, x_j \in X\setminus A$ con $x_i \neq x_j$, entonces $B_i \neq B_j$ (pues $x_i \notin B_j$) y por lo tanto $\{B_i\}_{i\in I}$ es una familia de abiertos no numerable, lo que contradice que $\cc{B}_n$ sea numerable pues estamos inyectando un conjunto no numerable en uno numerable. Luego $(X, T)$ no es 2AN.
        \end{itemize}
        \end{enumerate}
    \end{ejercicio}

    \begin{ejercicio}
    En $\mathbb{R}^2$ denotamos por $d$ a la distancia usual y definimos la familia de subconjuntos $T$ de la siguiente forma:
    \[
        T = \{U \subset \mathbb{R}^2 : (0,0) \notin U\} \cup \{U \subset \mathbb{R}^2 : \exists \varepsilon > 0 \text{ con } B_d((0,0), \varepsilon) \subset U\}.
    \]
    \begin{enumerate}
        \item (1,5 punto) Demuestra que $T$ es una topología en $\mathbb{R}^2$. Compárala con la topología usual de $\mathbb{R}^2$.\\
        Para ello vamos a probar qye $T\subseteq \cc{P}(\bb{R}^2)$ cumple las siguientes propiedades:
        \begin{enumerate}
            \item[(A1)] $\varnothing, \bb{R}^2 \in T$\\
            Está claro que $\varnothing \in T$ porque $(0,0) \notin \varnothing$. Por otro lado $\bb{R}^2 \in T$ porque $B_d((0,0), 1) \subset \bb{R}^2$.
            \item[(A2)] Si $\{U_i\}_{i\in I} \subset T$, entonces $\bigcup_{i\in I} U_i \in T$\\
            Sea $\{U_i\}_{i\in I} \subset T$, entonces existen dos casos:
            \begin{itemize}
                \item Si $(0,0) \notin U_i$ para todo $i\in I$, entonces $(0,0) \notin \bigcup_{i\in I} U_i$ y por lo tanto $\bigcup_{i\in I} U_i \in T$.
                \item Si existe un $i_0\in I$ tal que $(0,0) \in U_{i_0}$, entonces $\exists \varepsilon > 0$ tal que $B_d((0,0), \varepsilon) \subset U_{i_0}$. Por lo tanto $B_d((0,0), \varepsilon) \subset U_{i_0} \subset \bigcup_{i\in I} U_i$ y por lo tanto $\bigcup_{i\in I} U_i \in T$.
            \end{itemize}
            \item[(A3)] Si $U_1, U_2 \in T$, entonces $U_1 \cap U_2 \in T$\\
            Sea $U_1, U_2 \in T$. Si $U_1\cap U_2 = \varnothing \in T$ y ya lo tendríamos. Si $U_1\cap U_2 \neq \varnothing$, entonces distinguimos casos:
            \begin{itemize}
                \item Si $(0,0)\notin U_1$ (y/o análogamente $(0,0)\notin U_2$), entonces $(0,0) \notin U_1\cap U_2$ y por lo tanto $U_1 \cap U_2 \in T$.
                \item Si $(0,0) \in U_1, U_2$ entonces $\exists \delta, \epsilon > 0$ tales que $B_d((0,0), \delta) \subset U_1$ y $B_d((0,0), \epsilon) \subset U_2$. Sea $\varepsilon = \min\{\delta, \epsilon\}$, entonces $B_d((0,0), \varepsilon) \subset B_d((0,0), \delta) \cap B_d((0,0), \epsilon) \subset U_1 \cap U_2$ y por lo tanto $U_1 \cap U_2 \in T$.
            \end{itemize}
        \end{enumerate}
        Además es claro que $T_u\leq T$. Comprobémoslo de la siguiente manera: sea $B_1 = B_d((x,y), r)\in B_u$ con $(x,y) \in B_1$, $r>0$ entonces:
        \begin{itemize}
            \item Si $(x,y) = (0,0)$ entonces está claro que existe un $B_d((0,0), r') \in T$ con $r>r'>0$ tal que $(x,y)\in B_d((0,0), r') \subset B_d((0,0), r)$
            \item Si $(x,y) \neq (0,0)$ entonces para $d((x,y),(0,0)) > \epsilon > 0$ y $r > \epsilon$ se tiene $(0,0)\notin B_d((x,y), \epsilon) \Rightarrow B_d((x,y), \epsilon) \in T$ y por lo tanto $(x,y)\in B_d((x,y), \epsilon) \subset B_d((x,y), r)$
        \end{itemize}
        Por inclusión de las bases, $T_u\leq T$. Pero la otra implicación no es cierta, pues por ejemplo$\overline{B_d}((0,0), 1) \in T$ pero $\overline{B_d}((0,0), 1)\notin T_u$.
        
        \item (1 punto) Calcula una base de entornos de $(x,y) \in \mathbb{R}^2$ en $(\mathbb{R}^2, T)$:\\
        Sea $(x,y)\in\bb{R}^2$ entonces:
        \[\begin{cases}
            \cc{B}_{(x,y)} = \{\{(x,y)\}\} & \text{ si }(x,y)\neq (0,0)\\
            \cc{B}_{(0,0)} = \{B_d((0,0), r) : r > 0\}
        \end{cases}\]
        Vamos a comprobar que, efectivamente, es una base de entornos distinguiendo dos casos:
        \begin{itemize}
            \item $(x,y) = (0,0)$: Todo $B_d((0,0), r)$ es abierto en $T$ y contiene a $(0,0)$. Además, si $U$ es un entorno de $(0,0)$, existe un abierto $W\in T$ con $(0,0)\in W\subset U$, y por definición de $T$, $\exists \varepsilon > 0$ tal que $(0,0)\in V = B_d((0,0), \varepsilon) \subset W \subset U$ con $V\in\cc{B}_{(0,0)}$.
            \item $(x,y) \neq (0,0)$: Como $(0,0)\notin \{(x,y)\}$, se tiene que $V = \{(x,y)\}\in T$ es un entorno abierto de $(x,y)$. Así, para cualquier entorno $U$ de $(x,y)$ se cumple $(x,y)\in V\subset U$ con $V\in\cc{B}_{(x,y)}$.
        \end{itemize}
        Efectivamente $\cc{B}_{(x,y)}$ es una base de entornos de $(x,y)$ en $(\mathbb{R}^2, T)$.

        \item (1 punto) Calcula la adherencia y el interior de $\mathbb{Q} \times \{0\}$ y $\mathbb{Q} \times (0,1)$ en $(\mathbb{R}^2, T)$.
        \begin{itemize}
        \item $A = \mathbb{Q} \times \{0\}$\\
        Demostramos que $A^{\circ} = (\mathbb{Q}\setminus \{0\}) \times \{0\}$. Lo probaremos por doble implicación:
        \begin{itemize}
            \item[$\subseteq$] Basta ver que $(0,0) \notin A^{\circ}$. Si $(0,0) \in A^{\circ}$, existiría un abierto $U \in T$ tal que $(0,0) \in U \subset A$. Por definición de $T$, existiría $\varepsilon > 0$ tal que $B_d((0,0), \varepsilon) \subset U \subset A = \mathbb{Q} \times \{0\}$, lo cual es absurdo porque el punto $(0, \varepsilon/2) \in B_d((0,0), \varepsilon)$ no pertenece a $A$.
            \item[$\supseteq$] Como $(0,0) \notin (\mathbb{Q}\setminus \{0\}) \times \{0\}$, por definición de $T$ este conjunto es abierto en $(\mathbb{R}^2, T)$. Al estar contenido en $A$, se cumple $(\mathbb{Q}\setminus \{0\}) \times \{0\} \subset A^{\circ}$.
        \end{itemize}

        Por otro lado veamos que $\overline{A} = A = \mathbb{Q} \times \{0\}$ probando que $A$ es cerrado. Tenemos que $(0,0) \in A$, por lo que $(0,0) \notin \mathbb{R}^2 \setminus A$. Por definición de $T$, todo subconjunto que no contiene al origen es abierto, luego $\mathbb{R}^2 \setminus A \in T$ y por tanto $A$ es cerrado y $\overline{A} = A$.
        \begin{figure}[h!]
        \centering
        \begin{tikzpicture}[scale=1.6]
            % Ejes coordenados
            \draw[->, gray!70] (-1.6,0) -- (1.7,0) node[right, black] {$x$};
            \draw[->, gray!70] (0,-0.8) -- (0,1.2) node[above, black] {$y$};

            % Entorno básico en el origen: bola B_d((0,0), \varepsilon)
            \fill[red!15, opacity=0.7] (0,0) circle (0.55);
            \draw[red, dashed, thick] (0,0) circle (0.55);
            \node[red, font=\scriptsize, anchor=south west] at (-1.35, 0.42) {$B_d((0,0),\varepsilon)$};

            % Conjunto A = Q x {0} (representado como puntos densos en el eje X)
            \foreach \x in {-1.4, -1.3, ..., 1.5} {
                \fill[blue!80!black] (\x, 0) circle (0.7pt);
            }
            \node[blue!80!black, font=\small, anchor=north west] at (0.85, -0.55) {$A = \mathbb{Q} \times \{0\}$};

            % Punto origen (en A, pero no en A^\circ)
            \fill[blue!80!black] (0,0) circle (1.4pt);
            \node[font=\scriptsize, anchor=north east] at (-0.02, -0.02) {$(0,0)$};

            % Entorno básico fuera del origen: {(q,0)} aislado
            \fill[teal!80!black] (1.0, 0) circle (1.5pt);
            \draw[teal!80!black, thick] (1.0, 0) circle (3pt);
            \node[teal!80!black, font=\scriptsize, anchor=south] at (1.0, 0.1) {$\{(q,0)\}$};
        \end{tikzpicture}
        \caption{Representación de $A = \mathbb{Q} \times \{0\}$ y sus entornos críticos.}
        \end{figure}

        \item $B = \mathbb{Q} \times (0,1)$\\
        Es fácil ver que $B^{\circ} = B = \mathbb{Q} \times (0,1)$ pues $B\in T$. Para probarlo veamos que $(0,0)\notin B$, lo cual es inmediato ya que la segunda coordenada de $(0,0)$ es $0 \notin (0,1)$.

        Ahora veamos que $\overline{B} = (\mathbb{Q} \times (0,1)) \cup \{(0,0)\}$ por doble implicación:
        \begin{itemize}
            \item[$\supseteq$] Claramente $B \subset \overline{B}$. Para ver que $(0,0) \in \overline{B}$, sea $U \in T$ un entorno abierto cualquiera con $(0,0) \in U$. Por definición de $T$, existe $\varepsilon > 0$ tal que $B_d((0,0), \varepsilon) \subset U$. Tomando $r = \min\{\varepsilon/2, 1/2\} > 0$, se cumple $(0, r) \in B_d((0,0), \varepsilon) \subset U$ y además $(0, r) \in \mathbb{Q} \times (0,1) = B$. Luego $U \cap B \neq \varnothing$ y por tanto $(0,0) \in \overline{B}$.
            \item[$\subseteq$] Sea $F = B \cup \{(0,0)\}$. Su complementario $\mathbb{R}^2 \setminus F$ no contiene al $(0,0)$, por lo que $\mathbb{R}^2 \setminus F \in T$ y $F$ es cerrado. Como $B \subset F$ y $\overline{B}$ es el menor cerrado que contiene a $B$, se concluye que $\overline{B} \subset F$.
        \end{itemize}
        \end{itemize}

        \begin{figure}[h!]
        \centering
        \begin{tikzpicture}[scale=1.6]
            % Ejes coordenados
            \draw[->, gray!70] (-1.6,0) -- (1.7,0) node[right, black] {$x$};
            \draw[->, gray!70] (0,-0.7) -- (0,1.4) node[above, black] {$y$};

            % Entorno básico en el origen: bola B_d((0,0), \varepsilon)
            \fill[red!15, opacity=0.7] (0,0) circle (0.55);
            \draw[red, dashed, thick] (0,0) circle (0.55);
            \node[red, font=\scriptsize, anchor=north east] at (-0.35, -0.4) {$B_d((0,0),\varepsilon)$};

            % Conjunto B = Q x (0,1) (segmentos verticales abiertos en abscisas racionales)
            \foreach \x in {-1.4, -1.3, ..., 1.5} {
                \draw[blue!75!black, line width=0.6pt] (\x, 0) -- (\x, 1);
                \fill[white] (\x, 0) circle (0.7pt);
                \draw[blue!75!black, line width=0.4pt] (\x, 0) circle (0.7pt);
                \fill[white] (\x, 1) circle (0.7pt);
                \draw[blue!75!black, line width=0.4pt] (\x, 1) circle (0.7pt);
            }
            \node[blue!80!black, font=\small, anchor=west] at (1.52, 0.6) {$B = \mathbb{Q} \times (0,1)$};

            % Origen (en \overline{B}, fuera de B)
            \fill[red] (0,0) circle (1.4pt);

            % Entorno fuera del origen en (q,y) \in B: {(q,y)} \subset B
            \fill[orange!85!black] (-0.9, 0.65) circle (1.5pt);
            \draw[orange!85!black, thick] (-0.9, 0.65) circle (3pt);
            \node[orange!85!black, font=\scriptsize, anchor=east] at (-0.98, 0.65) {$\{(q,y)\}$};
        \end{tikzpicture}
        \caption{Representación de $B = \mathbb{Q} \times (0,1)$: $(0,0) \in \overline{B}$ porque toda bola $B_d((0,0),\varepsilon)$ corta a $B$ en $(0,r)$, pero los puntos $(q,0)$ con $q \neq 0$ no están en $\overline{B}$ al ser $V = \{(q,0)\}$ un entorno abierto disjunto de $B$.}
        \end{figure}
    \end{enumerate}
    \end{ejercicio}
    
\end{document}